\documentclass[11pt]{article}

\usepackage[margin=1in]{geometry}
\usepackage{amsmath, amssymb, amsthm}
\usepackage{hyperref}

\newtheorem*{takeaway}{Takeaway}
\newtheorem{theorem}{Theorem}
\newtheorem{corollary}{Corollary}
\newtheorem{proposition}{Proposition}

\title{Building the Irreducible-Error Theorem: A Guided Derivation}
\author{Cate Merfeld Dunham}
\date{\today}

\begin{document}
\maketitle

\section*{Overview}

This is a running log of a self-guided derivation exercise: working out, from
first principles and via guided questions rather than being handed the
answer, why a sensor that sometimes produces fully ambiguous readings caps
the best possible classification accuracy at a fixed ceiling. It is a
companion to, and independent re-derivation of, the result already written
up formally in the theorem note in this directory. This log is updated as
the exercise progresses, so it may contain open questions not yet resolved
at any given point.

\section{Warm-up: reasoning about a single reading}

\textbf{Setup considered.} Suppose you receive a single sensor reading, and
you are told that, given this specific reading, there is a 70\% chance the
true chemical is A and a 30\% chance it is B. You must guess A or B for this
one reading.

\textbf{Question asked.} Which would you guess, what is your probability of
being right, and can any strategy do better?

\textbf{Answer given.} Guess A, since 70\% is better than 30\%. No strategy
can beat the 70\% probability of being right.

\begin{takeaway}
For a single reading with known posterior chance $q$ of being class A, the
best possible probability of guessing correctly is $\max(q, 1-q)$, achieved
by guessing whichever class is more likely. No strategy can beat this for
that one reading.
\end{takeaway}

\section{Generalizing to a formula}

\textbf{Definition introduced.} For a general reading value $x$, let
\[
  p(x) = P(Y = A \mid X = x)
\]
be the probability that the true class of a sample is $A$, given that the
observed reading was $x$. (Confirmed explicitly: $Y$ is the true class label
and $X$ is the observed reading -- different kinds of objects, so
expressions like ``$X=Y$'' do not make sense; $p(x)$ is the conditional
probability $P(Y=A\mid X=x)$.)

\textbf{Question asked.} Generalizing the warm-up, write an expression for
the best possible accuracy at reading $x$, in terms of $p(x)$.

\textbf{Answer given.}
\[
  \max\big(p(x),\, 1-p(x)\big)
\]

\begin{takeaway}
The best possible accuracy at any single reading $x$ is $\max(p(x), 1-p(x))$.
\end{takeaway}

\section{Combining across all readings}

\textbf{Toy example (first version).} Readings take one of two appearances,
$x_1$ or $x_2$, each occurring 50\% of the time; best accuracy at $x_1$ is
90\%, at $x_2$ is 60\%. Question: out of 100 readings, how many correct
guesses overall, and what fraction?

Working through this example surfaced two important clarifications:

\begin{itemize}
  \item The accuracy numbers at $x_1$ and at $x_2$ are independent of each
  other -- each comes from its own pair of complementary probabilities
  ($p(x_1)$ and $1-p(x_1)$ sum to 1; separately, $p(x_2)$ and $1-p(x_2)$ sum
  to 1), but there is no required relationship between the $x_1$ pair and
  the $x_2$ pair.
  \item How often each reading \emph{appearance} occurs ($P(X=x_1)$,
  $P(X=x_2)$) is a different quantity from the overall prior on the true
  class ($P(Y=A)$, $P(Y=B)$), and the two need not match.
\end{itemize}

\textbf{Toy example (consistency-checked version).} Requiring $P(Y=A) =
P(Y=B) = 50\%$ constrains how the reading-appearance frequencies and the
per-reading posteriors can relate. One consistent instance: $P(X=x_1) =
0.2$, $P(X=x_2) = 0.8$, $p(x_1) = 0.9$, $p(x_2) = 0.4$, which averages back
out to $P(Y=A) = 0.2(0.9) + 0.8(0.4) = 0.5$ as required, giving the same
90\%/60\% accuracy pair as before but now on consistent footing.

\textbf{Further constraint added.} Additionally requiring $P(X=x_1) =
P(X=x_2) = 50\%$ forces
\[
  0.5\,p(x_1) + 0.5\,p(x_2) = 0.5 \quad\Longrightarrow\quad p(x_2) = 1 - p(x_1),
\]
i.e.\ the two per-reading posteriors become complements of each other. With
$p(x_1) = 0.9$, this forces $p(x_2) = 0.1$, and the best accuracy at $x_2$
becomes $\max(0.1, 0.9) = 90\%$ -- equal to the accuracy at $x_1$, rather
than a different value.

\textbf{Question asked.} Under this symmetric version (both readings equally
likely, both with 90\% best accuracy), what is the overall expected
accuracy out of 100 readings?

\textbf{Answer given.}
\[
  0.9 \times 0.5 + 0.9 \times 0.5 = 0.9 \quad (90\%)
\]

\begin{takeaway}
When equal priors ($P(Y=A)=P(Y=B)$) are combined with equal
reading-appearance frequencies, the per-reading posteriors are forced to be
complements of one another, which forces the best-possible accuracy to be
\emph{the same} at every reading appearance. In that case the overall
accuracy equals that common value regardless of how the reading-appearance
weights are distributed. This symmetry is expected to matter for the real
(sensor) problem.
\end{takeaway}

\textbf{Terminology introduced.} The operation of combining a per-reading
quantity across all possible readings, weighted by how likely each reading
is, is called the \emph{expected value} (or \emph{expectation}), written
$E_X[\cdot]$. The computation above is an instance of
\[
  E_X\big[\max(p(X), 1-p(X))\big].
\]

\section{The general inequality (Bayes-optimality bound)}

\textbf{Question asked.} Using the per-reading fact from
Section~2 (no strategy beats $\max(p(x),1-p(x))$ at a given $x$), write the
relationship between $P(g(X)=Y)$, for an arbitrary classifier $g$, and
$E_X[\max(p(X), 1-p(X))]$.

\textbf{Answer given.}
\[
  P(g(X) = Y) \;\le\; E_X\big[\max(p(X), 1-p(X))\big]
\]

\begin{takeaway}
For \emph{any} classifier $g$ (any function from readings to a predicted
class), its overall accuracy cannot exceed
$E_X[\max(p(X), 1-p(X))]$. This holds regardless of the corruption/ambiguity
model chosen -- it is a general fact about binary classification with a
known posterior $p(x) = P(Y=A\mid X=x)$.
\end{takeaway}

\section{Why the inequality holds}

\textbf{Sub-question asked.} At a specific reading $x$, if $g(x) = A$, what
is $P(\text{correct} \mid X=x)$ in terms of $p(x)$? What if $g(x) = B$
instead?

\textbf{Answer worked out.} If $g(x)=A$, the guess is correct exactly when
the true class really is $A$, so $P(\text{correct}\mid X=x) = P(Y=A\mid
X=x) = p(x)$. If $g(x) = B$, correctness requires the true class to be $B$,
so $P(\text{correct}\mid X=x) = 1-p(x)$. Either way, this per-reading
correctness probability is one of $\{p(x), 1-p(x)\}$, both of which are
$\le \max(p(x), 1-p(x))$ by definition.

\textbf{General averaging principle established.} Using a simpler warm-up
(if every number in a list is at most 10, the average of the list is at
most 10): if every term being averaged is bounded above by some ceiling
(possibly a different ceiling for each term), the average of the terms is
bounded above by the average of the ceilings.

\textbf{Combining.} The overall accuracy $P(g(X)=Y)$ is exactly the
expectation, over all readings, of the per-reading correctness probability:
\[
  P(g(X)=Y) = E_X\big[P(\text{correct}\mid X)\big].
\]
Since each term $P(\text{correct}\mid X=x) \le \max(p(x), 1-p(x))$, the
averaging principle gives
\[
  P(g(X)=Y) \;\le\; E_X\big[\max(p(X), 1-p(X))\big],
\]
which is exactly the inequality from Section~4 -- now understood as a
consequence of (a) a hard per-reading ceiling and (b) overall accuracy being
an average across readings.

\section{Achievability: the Bayes classifier}

\textbf{Question asked.} At reading $x$, $g(x)$ only has two possible
choices ($A$ or $B$), giving correctness probability $p(x)$ or $1-p(x)$
respectively. Is there always a choice matching $\max(p(x),1-p(x))$
exactly? Is anything stopping that choice from being made independently at
every $x$, as one single function $g$?

\textbf{Answer given.} Yes, there is always a choice that makes the
correctness probability equal to the ceiling; no, nothing prevents making
that choice independently at each reading.

\textbf{Explicit rule constructed.}
\[
  g^\ast(x) = \begin{cases} A & \text{if } p(x) \ge 1-p(x) \\ B & \text{otherwise.} \end{cases}
\]
(Here $g$, unstarred, denotes an arbitrary classifier -- the inequality
above holds for every such $g$ -- while $g^\ast$ denotes this specific,
deliberately constructed classifier.)

\textbf{Edge case checked.} At a reading where $p(x) = 1-p(x) = 0.5$ (a
tie), $\max(p(x),1-p(x)) = 0.5$ regardless of which class is guessed, so it
does not matter which branch a tie falls into; the $\ge$ in the rule above
resolves ties arbitrarily (into the $A$ branch) without affecting whether
the ceiling is met.

\section{The bound is tight}

\textbf{Question asked.} Given that $g^\ast$ meets $\max(p(x), 1-p(x))$
exactly at every reading $x$ (ties included): what does this establish
about the inequality $P(g(X)=Y) \le E_X[\max(p(X),1-p(X))]$ -- is
$E_X[\max(p(X),1-p(X))]$ merely \emph{a} valid upper bound, or is it the
best possible (tight) one?

\textbf{Answer given.} It can actually be reached.

\begin{takeaway}
Combining Section~5 ($P(g(X)=Y) \le E_X[\max(p(X),1-p(X))]$ for every
classifier $g$) with Section~6 ($g^\ast$ achieves this value exactly)
establishes that $E_X[\max(p(X),1-p(X))]$ is not merely a valid ceiling but
the tight maximum: the largest accuracy any classifier can achieve, reached
with equality by $g^\ast$. This general fact -- true for any binary
classification problem with posterior $p(x)=P(Y=A\mid X=x)$, independent of
any corruption or ambiguity model -- is the statement usually called
\emph{Bayes-optimality of the MAP classifier}.
\end{takeaway}

\section{Part II: is the ``smudge'' provable in general?}

Having finished the general Bayes-optimality result above, the discussion
turned to a second, separate question raised by the PCA figure: must the
gradient of points between the two target clusters (the ``smudge'') always
be present, and can that be shown in a way that applies to any predictor,
not just the specific MLP encoder from the experiment? The answer is yes,
but it requires a shift in framing.

\textbf{Key framing shift.} Everything proved in Part I concerns a
\emph{discrete decision} $g(x)\in\{A,B\}$, scored by whether it is exactly
right (0--1 loss / accuracy). ``Smudge'' only makes sense for a
\emph{continuous-valued} output landing anywhere between two points --
which is actually what the real encoder produces (a continuous vector,
trained by minimizing \emph{squared error}, not by making a discrete
decision). So Part II re-derives the analogous ``best possible output''
fact, but for squared-error loss instead of 0--1 loss, using the same
pointwise-optimize-then-average proof technique as Part I.

\subsection{Warm-up: a single reading, squared-error loss}

\textbf{Setup.} As in the very first warm-up, suppose at a specific reading
$x$, the target $T$ equals $1$ (class $A$) with probability 70\%, and $0$
(class $B$) with probability 30\%. This time, instead of guessing $A$ or
$B$ to maximize the probability of being exactly right, the task is to pick
a single real number $v$ minimizing the \emph{expected squared error}
$E[(v-T)^2]$.

\textbf{Question asked.} What value of $v$ minimizes $E[(v-T)^2]$, and does
the old ``commit to the more likely outcome'' reasoning (which gave $v=1$
in the classification version) still apply?

\textbf{Answer given.} $v = 0.7$ -- described initially as ``70\% confident
it's 1, because that's true 70\% of the time,'' then sharpened to: this is
the \emph{expected value} of $T$, $E[T] = 1(0.7) + 0(0.3) = 0.7$, which is a
different operation from ``pick the more likely single outcome'' (that
would give $v=1$, not $0.7$).

\textbf{Terminology reviewed.} The \emph{expectation} of a random variable
$Z$ taking values $z_1, z_2, \ldots$ with probabilities $P(Z=z_1),
P(Z=z_2),\ldots$ is $E[Z] = \sum_i z_i P(Z=z_i)$ -- the same operation used
earlier for $E_X[\max(p(X),1-p(X))]$, now applied to $T$ instead of to a
per-reading accuracy value.

\begin{takeaway}
For a single reading with a binary target $T$, the squared-error-minimizing
output is $E[T]$, the probability-weighted average of $T$'s possible
values -- \emph{not} $\max$ of the two probabilities as in the
classification case. These are genuinely different optimization problems
with genuinely different optimal answers.
\end{takeaway}

\subsection{General formula}

\textbf{Question asked.} For a general reading $x$ where $T=1$ with
probability $p(x)$ and $T=0$ with probability $1-p(x)$, write the general
formula for the optimal $v$.

\textbf{Answer given (general form).}
\[
  E[T] = t_1\, p(x) + t_0\,(1-p(x))
\]

\textbf{Specialized.} Substituting the actual values $t_1=1$, $t_0=0$:
\[
  v^\ast(x) = E[T] = p(x).
\]

\begin{takeaway}
The squared-error-optimal output at reading $x$ is simply $p(x)$ itself --
a strikingly simple formula, and one that is central to the smudge
question (see the open question below).
\end{takeaway}

\section{When must the smudge appear?}

\textbf{Question asked.} Given $v^\ast(x) = p(x)$: when does this equal one
of the two ``pure'' target values ($0$ or $1$), and when is it forced to
sit strictly between them?

\textbf{Answer given.} If $p(x)$ is $1$ (or $0$), $v^\ast(x)$ would be $1$
(or $0$). If $p(x)$ is between $0$ and $1$, $v^\ast(x)$ will be between
them. So if there is any ambiguity, the optimal output must be smudged.

\begin{takeaway}
$v^\ast(x) \in \{0,1\}$ iff $p(x) \in \{0,1\}$ (no ambiguity); whenever
$0<p(x)<1$, $v^\ast(x)$ is strictly between $0$ and $1$. The smudge is
mathematically necessary exactly where genuine ambiguity exists.
\end{takeaway}

\section{Extending to the full one-hot vector}

\textbf{Question asked.} With $T=e_A$ (probability $p(x)$) or $T=e_B$
(probability $1-p(x)$), write $v^\ast(x)$ as a combination of $p(x)$,
$e_A$, and $e_B$.

\textbf{Answer given.}
\[
  v^\ast(x) = e_A\, p(x) + e_B\,(1-p(x)).
\]

\begin{takeaway}
The vector-valued optimum is a convex combination of the two one-hot
targets, i.e.\ a point on the line segment $[e_A, e_B]$, weighted by
$p(x)$. Restricted to just the two coordinates corresponding to $A$ and
$B$ (all others are always $0$, since $T$ never has mass there), this is
$(p(x), 1-p(x))$ -- a point that coincides with an endpoint of the segment
only when $p(x)\in\{0,1\}$, and otherwise sits strictly in its interior.
\end{takeaway}

\section{Formal statement}

\begin{theorem}[Bayes-optimality of the MAP classifier]
Let $X$ and $Y$ be jointly distributed random variables with $Y\in\{A,B\}$,
and define $p(x) := P(Y=A\mid X=x)$. Then for every measurable function
$g:\mathcal X \to \{A,B\}$,
\[
  P\big(g(X)=Y\big) \;\le\; \mathbb{E}_X\big[\max(p(X), 1-p(X))\big],
\]
and equality is attained by
\[
  g^\ast(x) = \begin{cases} A & \text{if } p(x) \ge 1-p(x) \\ B & \text{otherwise.} \end{cases}
\]
\end{theorem}

\begin{proof}
We prove the inequality first, then achievability.

\textbf{The inequality.} Fix a measurable $g:\mathcal X\to\{A,B\}$ and a
reading $x$. Since $g(x)$ is either $A$ or $B$, exactly one of two cases
holds: if $g(x)=A$, then $g$ is correct at $x$ precisely when $Y=A$, giving
$P(g(X)=Y\mid X=x)=p(x)$; if $g(x)=B$, correctness requires $Y=B$, giving
$P(g(X)=Y\mid X=x)=1-p(x)$. In either case the correctness probability
equals one of $p(x)$ or $1-p(x)$, so
\[
  P(g(X)=Y\mid X=x) \;\le\; \max(p(x), 1-p(x)).
\]
Taking expectations over $X$ and applying the law of total probability,
$P(g(X)=Y) = \mathbb{E}_X[P(g(X)=Y\mid X)]$, gives
\[
  P(g(X)=Y) \;\le\; \mathbb{E}_X\big[\max(p(X), 1-p(X))\big].
\]
Since $g$ was an arbitrary measurable function, this holds for every
classifier.

\textbf{Achievability.} Define $g^\ast(x) = A$ if $p(x)\ge 1-p(x)$, and
$g^\ast(x)=B$ otherwise -- i.e.\ $g^\ast$ always selects whichever of
$p(x)$, $1-p(x)$ is larger. By the case analysis above, this means
$P(g^\ast(X)=Y\mid X=x) = \max(p(x),1-p(x))$ for every $x$: not merely
bounded by the per-reading ceiling, but equal to it. Taking expectations,
\[
  P(g^\ast(X)=Y) = \mathbb{E}_X\big[\max(p(X), 1-p(X))\big],
\]
so $g^\ast$ attains the bound with equality, and the inequality above is
tight.
\end{proof}

\begin{theorem}[$L^2$-optimality of the conditional expectation]
Let $T$ be a $\{e_A, e_B\}$-valued random vector ($e_A, e_B \in
\mathbb{R}^n$ the one-hot targets) with $P(T=e_A\mid X=x) = p(x)$ and
$P(T=e_B\mid X=x) = 1-p(x)$. Then, over all measurable $f:\mathcal X \to
\mathbb{R}^n$, $\mathbb{E}\big[\|f(X)-T\|^2\big]$ is minimized (a.e.\
uniquely) by
\[
  f^\ast(x) = \mathbb{E}[T\mid X=x] = p(x)\,e_A + (1-p(x))\,e_B.
\]
\end{theorem}

\begin{proof}[Proof of Theorem 2]
\textbf{Step 1: one fixed reading.} Fix a reading $x$ and a candidate output
$v\in\mathbb{R}^n$. Given $X=x$, $T$ equals $e_A$ with probability $p(x)$ and
$e_B$ with probability $1-p(x)$, so
\[
  \mathbb{E}\big[\|v-T\|^2 \mid X=x\big]
  = p(x)\,\|v-e_A\|^2 + (1-p(x))\,\|v-e_B\|^2.
\]
Expand each term using $\|v-e\|^2 = \|v\|^2 - 2\,v\cdot e + \|e\|^2$, and
write $m := p(x)\,e_A + (1-p(x))\,e_B$. The terms involving $v$ combine to
\[
  \mathbb{E}\big[\|v-T\|^2 \mid X=x\big] = \|v\|^2 - 2\,v\cdot m + C,
  \qquad C := p(x)\|e_A\|^2 + (1-p(x))\|e_B\|^2,
\]
where $C$ does not depend on $v$. Completing the square, since
$\|v-m\|^2 = \|v\|^2 - 2\,v\cdot m + \|m\|^2$,
\[
  \mathbb{E}\big[\|v-T\|^2 \mid X=x\big] = \|v-m\|^2 + \big(C - \|m\|^2\big).
\]
The second term does not depend on $v$, and the first is nonnegative and
equals zero exactly when $v=m$. So $v=m=f^\ast(x)$ is the unique minimizer at
this reading.

(Scalar warm-up of the same algebra, with $T\in\{0,1\}$ and $p(x)=0.7$:
$0.7(v-1)^2 + 0.3v^2 = (v-0.7)^2 + 0.21$, minimized at $v=0.7$, with
unavoidable error $0.21$.)

\textbf{Step 2: combine across readings.} By the law of total expectation,
for any measurable $f$,
\[
  \mathbb{E}\big[\|f(X)-T\|^2\big]
  = \mathbb{E}_X\Big[\mathbb{E}\big[\|f(X)-T\|^2 \mid X\big]\Big].
\]
By Step 1, for every reading $x$,
\[
  \mathbb{E}\big[\|f(x)-T\|^2 \mid X=x\big] \;\ge\;
  \mathbb{E}\big[\|f^\ast(x)-T\|^2 \mid X=x\big].
\]
Since an average of terms that are each bounded below is bounded below by
the average of those bounds, taking $\mathbb{E}_X$ gives
\[
  \mathbb{E}\big[\|f(X)-T\|^2\big] \;\ge\; \mathbb{E}\big[\|f^\ast(X)-T\|^2\big]
\]
for every measurable $f$, so $f^\ast$ minimizes the overall expected squared
error.

\textbf{Uniqueness.} Suppose some $f$ attains the same minimum. Define the
per-reading gap
\[
  D(x) := \mathbb{E}\big[\|f(x)-T\|^2 \mid X=x\big]
        - \mathbb{E}\big[\|f^\ast(x)-T\|^2 \mid X=x\big].
\]
By Step 1, $D(x)\ge 0$ for every $x$, and by assumption $\mathbb{E}_X[D(X)]=0$.
A nonnegative quantity whose average is zero must equal zero except possibly
on a set of readings of probability zero, so $D(x)=0$ for almost every $x$.
By Step 1, the only $v$ that achieves the minimum at a reading $x$ is
$m=f^\ast(x)$, so $f(x)=f^\ast(x)$ for almost every $x$. $\blacksquare$
\end{proof}

\begin{corollary}[Restricted to the two active coordinates]
Writing $f^\ast_A, f^\ast_B$ for the coordinates of $f^\ast$ corresponding
to $e_A, e_B$:
\[
  \big(f^\ast_A(x),\, f^\ast_B(x)\big) = \big(p(x),\, 1-p(x)\big) \in
  \mathrm{conv}\{(1,0),(0,1)\},
\]
which equals an extreme point of that segment iff $p(x)\in\{0,1\}$, and
lies in its relative interior whenever $0<p(x)<1$.
\end{corollary}

\begin{proof}[Proof of the Corollary]
By Theorem~2, $f^\ast(x) = p(x)\,e_A + (1-p(x))\,e_B$, and $e_A, e_B$ have a
$1$ in their own coordinate and $0$ in the other's, so the $A$ and $B$
coordinates of $f^\ast(x)$ are $p(x)$ and $1-p(x)$. This point equals
$(1,0)$ or $(0,1)$ exactly when $p(x)\in\{0,1\}$; for $0<p(x)<1$, both
coordinates are strictly between $0$ and $1$, so the point lies in the
interior of the segment.
\end{proof}

\begin{proposition}[Cost of not smudging]
In the setting of Theorem~2, for every measurable $f$,
\[
  \mathbb{E}\big[\|f(X)-T\|^2\big]
  = \mathbb{E}\big[\|f^\ast(X)-T\|^2\big]
  + \mathbb{E}\big[\|f(X)-f^\ast(X)\|^2\big].
\]
Consequently, if $f(x)\neq f^\ast(x)$ on a set of readings of positive
probability, then $\mathbb{E}\big[\|f(X)-T\|^2\big] >
\mathbb{E}\big[\|f^\ast(X)-T\|^2\big]$. In particular, a predictor that
outputs a pure corner at readings where $0<p(x)<1$ has strictly larger
expected squared error than $f^\ast$.
\end{proposition}

\begin{proof}
Fix a reading $x$. Step~1 of the proof of Theorem~2 gives, for any vector
$v$,
\[
  \mathbb{E}\big[\|v-T\|^2 \mid X=x\big] = \|v-f^\ast(x)\|^2 + \big(C-\|f^\ast(x)\|^2\big),
\]
where the second term does not depend on $v$. Taking $v=f(x)$ and $v=f^\ast(x)$
and subtracting,
\[
  \mathbb{E}\big[\|f(x)-T\|^2 \mid X=x\big] - \mathbb{E}\big[\|f^\ast(x)-T\|^2 \mid X=x\big]
  = \|f(x)-f^\ast(x)\|^2.
\]
Replace $x$ by the random reading $X$ and average both sides with $\mathbb{E}_X$.
By the law of total expectation,
$\mathbb{E}_X\big[\mathbb{E}[\|g(X)-T\|^2 \mid X]\big] = \mathbb{E}\big[\|g(X)-T\|^2\big]$
for $g=f$ and $g=f^\ast$, which gives
\[
  \mathbb{E}\big[\|f(X)-T\|^2\big] - \mathbb{E}\big[\|f^\ast(X)-T\|^2\big]
  = \mathbb{E}\big[\|f(X)-f^\ast(X)\|^2\big],
\]
the claimed identity. The last term is an average of nonnegative numbers, so it
is positive whenever $\|f(X)-f^\ast(X)\|^2>0$ on a set of positive
probability. If $f$ outputs a pure corner where $0<p(x)<1$, then $f(x)$ differs
from the strictly interior point $f^\ast(x)$ there, so $\|f(x)-f^\ast(x)\|^2>0$.
\end{proof}

\section{Empirical checks against the actual experiment}

A few points that came up while inspecting the real saved encoder outputs
(\texttt{mnist\_6\_8\_embeddings.pt}), relevant to the theorem above:

\begin{itemize}
  \item Restricting the PCA fit itself to dims 6 and 8, or better, plotting
  those two dimensions directly with no PCA at all, reproduces the
  corollary's geometry exactly: points fall almost exactly on the segment
  between $(1,0)$ and $(0,1)$ (measured deviation from
  $\text{dim}_6+\text{dim}_8=1$: mean $\approx 0$, max $0.024$ on train,
  $0.063$ on test).
  \item Points near the middle of that segment are not purely explained by
  label corruption: among the most ambiguous train points, $37.5\%$ were
  mislabeled versus a $19.8\%$ base rate -- corruption is overrepresented
  there, as expected, but several genuinely clean, correctly-assigned
  samples also land near the middle, reflecting ordinary visual ambiguity
  in handwriting rather than the corruption mechanism.
  \item The theorem's ``never exactly $0$ or $1$'' claim is about the true,
  infinite-precision optimum. In the actual saved \texttt{float32} tensors,
  thousands of entries are exactly $1.0$ (or $0.0$) -- a floating-point
  rounding artifact (the true sigmoid output can be closer to $1$ than the
  nearest representable float32 value other than $1.0$ itself), not a
  contradiction of the theorem.
  \item \textbf{Flipped-label test set} (same test images, the same $20\%$
  corruption process as training; this matches the unreliable-sensor scenario,
  where test data is as ambiguous as train data). Accuracy against the
  flipped labels was $71.4\%$, below the $1-q = 79.97\%$ ceiling of
  Theorem~1, as it must be. The model recovers the true digit on about
  $84\%$ of images, and $0.84(0.8)+0.16(0.2)\approx 0.70$ predicts the
  observed agreement. Mean squared error against the flipped targets was
  $0.0538$. If every reading had $p(x)=0.8$ and the model output
  $(0.8,0.2)$, the floor would be $2q(1-q)/10 = 0.032$; a predictor snapping
  perfectly to the true digit's corner would score $0.04$. The difference,
  $0.008 = \mathbb{E}\|f-f^\ast\|^2/10$, is the cost of not smudging from the
  proposition above.
\end{itemize}

\section{Scope decision}

The general results above are complete. Deriving the concrete $80\%$ ceiling
from a specific swap-corruption model (computing $p(x)$ explicitly for the
sensor scenario) was considered and deliberately dropped; the log stops at the
general theorems.

\end{document}
